\documentclass[12pt,a4paper]{article}
\usepackage[top=1.8cm,bottom=1.8cm,left=1.8cm,right=1.8cm]{geometry}
\usepackage{fontspec}
\setmainfont{TeX Gyre Termes}
\usepackage{amsmath,amssymb}
\usepackage{fancyhdr}
\usepackage{tasks}
\usepackage{tikz}
\settasks{label=\textbf{\Alph*.},label-width=1.5em,item-indent=2em}
\setlength{\parindent}{0pt}
\setlength{\parskip}{3pt}
\pagestyle{fancy}\fancyhf{}\renewcommand{\headrulewidth}{0pt}
\fancyhead[L]{\small Ngân hàng bài tập Toán 9 -- 9H2-C6-B20-D02}
\cfoot{\small Trang \thepage}
\newif\ifloigiai
\loigiaitrue
\newenvironment{cauhoi}[2]{%
  \par\medskip\noindent\rule{\textwidth}{0.4pt}\par\nopagebreak
  \noindent\textbf{[#1]}\hfill{\scriptsize\ttfamily #2}\par\nopagebreak\smallskip}{\par}
\newcommand{\loigiai}[1]{\ifloigiai\par\smallskip\noindent\textit{Lời giải.} #1\par\fi}
\newcommand{\dapan}[1]{\ifloigiai\par\smallskip\noindent\textit{Đáp án:} \textbf{#1}\par\fi}
\begin{document}
\begin{center}
{\Large\textbf{NGÂN HÀNG BÀI TẬP TOÁN 9}}\\[3pt]
\textbf{Chương 6. Hàm số y = ax² (a ≠ 0). Phương trình bậc hai một ẩn}\\[2pt]
\textbf{Bài 20. Định lí Viète và ứng dụng}\\[2pt]
\textbf{Dạng D02: Biểu thức đối xứng hai nghiệm}\\[2pt]
\textit{Mã dạng: 9H2-C6-B20-D02 \quad 50 bài tự luận + 50 câu trắc nghiệm \quad (VD: 50)}
\end{center}
\textbf{Quy ước mã câu:} \textit{khối, học kì, chương, bài, dạng, loại (TL/TN), mức độ, số thứ tự}.\par Ví dụ: \texttt{9H2-C6-B20-D02-TL-NB-001}.\par Để in bản cho học sinh, đặt \verb|\loigiaifalse| ngay sau \verb|\loigiaitrue|.
\section*{Phần I. Tự luận}
\begin{cauhoi}{Bài 1}{9H2-C6-B20-D02-TL-VD-001}
Gọi $x_1$, $x_2$ là hai nghiệm của phương trình $x^{2}+6x+8=0$. Tính $x_1^{2}+x_2^{2}$ (không giải phương trình).
\loigiai{Theo định lí Viète: $x_1+x_2=-6$, $x_1x_2=8$. Suy ra $x_1^{2}+x_2^{2}=(x_1+x_2)^{2}-2x_1x_2=-6^{2}-2\cdot(8)=20$.}
\end{cauhoi}
\begin{cauhoi}{Bài 2}{9H2-C6-B20-D02-TL-VD-002}
Gọi $x_1$, $x_2$ là hai nghiệm của phương trình $x^{2}-4x-5=0$. Tính $x_1^{2}+x_2^{2}$ (không giải phương trình).
\loigiai{Theo định lí Viète: $x_1+x_2=4$, $x_1x_2=-5$. Suy ra $x_1^{2}+x_2^{2}=(x_1+x_2)^{2}-2x_1x_2=4^{2}-2\cdot(-5)=26$.}
\end{cauhoi}
\begin{cauhoi}{Bài 3}{9H2-C6-B20-D02-TL-VD-003}
Gọi $x_1$, $x_2$ là hai nghiệm của phương trình $x^{2}-x-56=0$. Tính $x_1^{2}+x_2^{2}$ (không giải phương trình).
\loigiai{Theo định lí Viète: $x_1+x_2=1$, $x_1x_2=-56$. Suy ra $x_1^{2}+x_2^{2}=(x_1+x_2)^{2}-2x_1x_2=1^{2}-2\cdot(-56)=113$.}
\end{cauhoi}
\begin{cauhoi}{Bài 4}{9H2-C6-B20-D02-TL-VD-004}
Gọi $x_1$, $x_2$ là hai nghiệm của phương trình $x^{2}+9x+8=0$. Tính $x_1^{2}+x_2^{2}$ (không giải phương trình).
\loigiai{Theo định lí Viète: $x_1+x_2=-9$, $x_1x_2=8$. Suy ra $x_1^{2}+x_2^{2}=(x_1+x_2)^{2}-2x_1x_2=-9^{2}-2\cdot(8)=65$.}
\end{cauhoi}
\begin{cauhoi}{Bài 5}{9H2-C6-B20-D02-TL-VD-005}
Gọi $x_1$, $x_2$ là hai nghiệm của phương trình $x^{2}+7x=0$. Tính $x_1^{2}+x_2^{2}$ (không giải phương trình).
\loigiai{Theo định lí Viète: $x_1+x_2=-7$, $x_1x_2=0$. Suy ra $x_1^{2}+x_2^{2}=(x_1+x_2)^{2}-2x_1x_2=-7^{2}-2\cdot(0)=49$.}
\end{cauhoi}
\begin{cauhoi}{Bài 6}{9H2-C6-B20-D02-TL-VD-006}
Gọi $x_1$, $x_2$ là hai nghiệm của phương trình $x^{2}-4=0$. Tính $x_1^{2}+x_2^{2}$ (không giải phương trình).
\loigiai{Theo định lí Viète: $x_1+x_2=0$, $x_1x_2=-4$. Suy ra $x_1^{2}+x_2^{2}=(x_1+x_2)^{2}-2x_1x_2=0^{2}-2\cdot(-4)=8$.}
\end{cauhoi}
\begin{cauhoi}{Bài 7}{9H2-C6-B20-D02-TL-VD-007}
Gọi $x_1$, $x_2$ là hai nghiệm của phương trình $x^{2}-11x+24=0$. Tính $x_1^{2}+x_2^{2}$ (không giải phương trình).
\loigiai{Theo định lí Viète: $x_1+x_2=11$, $x_1x_2=24$. Suy ra $x_1^{2}+x_2^{2}=(x_1+x_2)^{2}-2x_1x_2=11^{2}-2\cdot(24)=73$.}
\end{cauhoi}
\begin{cauhoi}{Bài 8}{9H2-C6-B20-D02-TL-VD-008}
Gọi $x_1$, $x_2$ là hai nghiệm của phương trình $x^{2}-4x+3=0$. Tính $x_1^{2}+x_2^{2}$ (không giải phương trình).
\loigiai{Theo định lí Viète: $x_1+x_2=4$, $x_1x_2=3$. Suy ra $x_1^{2}+x_2^{2}=(x_1+x_2)^{2}-2x_1x_2=4^{2}-2\cdot(3)=10$.}
\end{cauhoi}
\begin{cauhoi}{Bài 9}{9H2-C6-B20-D02-TL-VD-009}
Gọi $x_1$, $x_2$ là hai nghiệm của phương trình $x^{2}-25=0$. Tính $x_1^{2}+x_2^{2}$ (không giải phương trình).
\loigiai{Theo định lí Viète: $x_1+x_2=0$, $x_1x_2=-25$. Suy ra $x_1^{2}+x_2^{2}=(x_1+x_2)^{2}-2x_1x_2=0^{2}-2\cdot(-25)=50$.}
\end{cauhoi}
\begin{cauhoi}{Bài 10}{9H2-C6-B20-D02-TL-VD-010}
Gọi $x_1$, $x_2$ là hai nghiệm của phương trình $x^{2}-8x=0$. Tính $x_1^{2}+x_2^{2}$ (không giải phương trình).
\loigiai{Theo định lí Viète: $x_1+x_2=8$, $x_1x_2=0$. Suy ra $x_1^{2}+x_2^{2}=(x_1+x_2)^{2}-2x_1x_2=8^{2}-2\cdot(0)=64$.}
\end{cauhoi}
\begin{cauhoi}{Bài 11}{9H2-C6-B20-D02-TL-VD-011}
Gọi $x_1$, $x_2$ là hai nghiệm của phương trình $x^{2}+4x-21=0$. Tính $x_1^{2}+x_2^{2}$ (không giải phương trình).
\loigiai{Theo định lí Viète: $x_1+x_2=-4$, $x_1x_2=-21$. Suy ra $x_1^{2}+x_2^{2}=(x_1+x_2)^{2}-2x_1x_2=-4^{2}-2\cdot(-21)=58$.}
\end{cauhoi}
\begin{cauhoi}{Bài 12}{9H2-C6-B20-D02-TL-VD-012}
Gọi $x_1$, $x_2$ là hai nghiệm của phương trình $x^{2}-3x-4=0$. Tính $x_1^{2}+x_2^{2}$ (không giải phương trình).
\loigiai{Theo định lí Viète: $x_1+x_2=3$, $x_1x_2=-4$. Suy ra $x_1^{2}+x_2^{2}=(x_1+x_2)^{2}-2x_1x_2=3^{2}-2\cdot(-4)=17$.}
\end{cauhoi}
\begin{cauhoi}{Bài 13}{9H2-C6-B20-D02-TL-VD-013}
Gọi $x_1$, $x_2$ là hai nghiệm của phương trình $x^{2}-12x+35=0$. Tính $x_1^{2}+x_2^{2}$ (không giải phương trình).
\loigiai{Theo định lí Viète: $x_1+x_2=12$, $x_1x_2=35$. Suy ra $x_1^{2}+x_2^{2}=(x_1+x_2)^{2}-2x_1x_2=12^{2}-2\cdot(35)=74$.}
\end{cauhoi}
\begin{cauhoi}{Bài 14}{9H2-C6-B20-D02-TL-VD-014}
Gọi $x_1$, $x_2$ là hai nghiệm của phương trình $x^{2}+8x=0$. Tính $x_1^{2}+x_2^{2}$ (không giải phương trình).
\loigiai{Theo định lí Viète: $x_1+x_2=-8$, $x_1x_2=0$. Suy ra $x_1^{2}+x_2^{2}=(x_1+x_2)^{2}-2x_1x_2=-8^{2}-2\cdot(0)=64$.}
\end{cauhoi}
\begin{cauhoi}{Bài 15}{9H2-C6-B20-D02-TL-VD-015}
Gọi $x_1$, $x_2$ là hai nghiệm của phương trình $x^{2}+x=0$. Tính $x_1^{2}+x_2^{2}$ (không giải phương trình).
\loigiai{Theo định lí Viète: $x_1+x_2=-1$, $x_1x_2=0$. Suy ra $x_1^{2}+x_2^{2}=(x_1+x_2)^{2}-2x_1x_2=-1^{2}-2\cdot(0)=1$.}
\end{cauhoi}
\begin{cauhoi}{Bài 16}{9H2-C6-B20-D02-TL-VD-016}
Gọi $x_1$, $x_2$ là hai nghiệm của phương trình $x^{2}+x-42=0$. Tính $x_1^{2}+x_2^{2}$ (không giải phương trình).
\loigiai{Theo định lí Viète: $x_1+x_2=-1$, $x_1x_2=-42$. Suy ra $x_1^{2}+x_2^{2}=(x_1+x_2)^{2}-2x_1x_2=-1^{2}-2\cdot(-42)=85$.}
\end{cauhoi}
\begin{cauhoi}{Bài 17}{9H2-C6-B20-D02-TL-VD-017}
Gọi $x_1$, $x_2$ là hai nghiệm của phương trình $x^{2}+2x-48=0$. Tính $x_1^{2}+x_2^{2}$ (không giải phương trình).
\loigiai{Theo định lí Viète: $x_1+x_2=-2$, $x_1x_2=-48$. Suy ra $x_1^{2}+x_2^{2}=(x_1+x_2)^{2}-2x_1x_2=-2^{2}-2\cdot(-48)=100$.}
\end{cauhoi}
\begin{cauhoi}{Bài 18}{9H2-C6-B20-D02-TL-VD-018}
Gọi $x_1$, $x_2$ là hai nghiệm của phương trình $x^{2}-x-30=0$. Tính $x_1^{2}+x_2^{2}$ (không giải phương trình).
\loigiai{Theo định lí Viète: $x_1+x_2=1$, $x_1x_2=-30$. Suy ra $x_1^{2}+x_2^{2}=(x_1+x_2)^{2}-2x_1x_2=1^{2}-2\cdot(-30)=61$.}
\end{cauhoi}
\begin{cauhoi}{Bài 19}{9H2-C6-B20-D02-TL-VD-019}
Gọi $x_1$, $x_2$ là hai nghiệm của phương trình $x^{2}+x-20=0$. Tính $x_1^{2}+x_2^{2}$ (không giải phương trình).
\loigiai{Theo định lí Viète: $x_1+x_2=-1$, $x_1x_2=-20$. Suy ra $x_1^{2}+x_2^{2}=(x_1+x_2)^{2}-2x_1x_2=-1^{2}-2\cdot(-20)=41$.}
\end{cauhoi}
\begin{cauhoi}{Bài 20}{9H2-C6-B20-D02-TL-VD-020}
Gọi $x_1$, $x_2$ là hai nghiệm của phương trình $x^{2}-2x-48=0$. Tính $x_1^{2}+x_2^{2}$ (không giải phương trình).
\loigiai{Theo định lí Viète: $x_1+x_2=2$, $x_1x_2=-48$. Suy ra $x_1^{2}+x_2^{2}=(x_1+x_2)^{2}-2x_1x_2=2^{2}-2\cdot(-48)=100$.}
\end{cauhoi}
\begin{cauhoi}{Bài 21}{9H2-C6-B20-D02-TL-VD-021}
Gọi $x_1$, $x_2$ là hai nghiệm của phương trình $x^{2}+13x+42=0$. Tính $x_1^{2}+x_2^{2}$ (không giải phương trình).
\loigiai{Theo định lí Viète: $x_1+x_2=-13$, $x_1x_2=42$. Suy ra $x_1^{2}+x_2^{2}=(x_1+x_2)^{2}-2x_1x_2=-13^{2}-2\cdot(42)=85$.}
\end{cauhoi}
\begin{cauhoi}{Bài 22}{9H2-C6-B20-D02-TL-VD-022}
Gọi $x_1$, $x_2$ là hai nghiệm của phương trình $x^{2}+4x-32=0$. Tính $x_1^{2}+x_2^{2}$ (không giải phương trình).
\loigiai{Theo định lí Viète: $x_1+x_2=-4$, $x_1x_2=-32$. Suy ra $x_1^{2}+x_2^{2}=(x_1+x_2)^{2}-2x_1x_2=-4^{2}-2\cdot(-32)=80$.}
\end{cauhoi}
\begin{cauhoi}{Bài 23}{9H2-C6-B20-D02-TL-VD-023}
Gọi $x_1$, $x_2$ là hai nghiệm của phương trình $x^{2}+4x-5=0$. Tính $x_1^{2}+x_2^{2}$ (không giải phương trình).
\loigiai{Theo định lí Viète: $x_1+x_2=-4$, $x_1x_2=-5$. Suy ra $x_1^{2}+x_2^{2}=(x_1+x_2)^{2}-2x_1x_2=-4^{2}-2\cdot(-5)=26$.}
\end{cauhoi}
\begin{cauhoi}{Bài 24}{9H2-C6-B20-D02-TL-VD-024}
Gọi $x_1$, $x_2$ là hai nghiệm của phương trình $x^{2}-3x-18=0$. Tính $x_1^{2}+x_2^{2}$ (không giải phương trình).
\loigiai{Theo định lí Viète: $x_1+x_2=3$, $x_1x_2=-18$. Suy ra $x_1^{2}+x_2^{2}=(x_1+x_2)^{2}-2x_1x_2=3^{2}-2\cdot(-18)=45$.}
\end{cauhoi}
\begin{cauhoi}{Bài 25}{9H2-C6-B20-D02-TL-VD-025}
Gọi $x_1$, $x_2$ là hai nghiệm của phương trình $x^{2}-6x-16=0$. Tính $x_1^{2}+x_2^{2}$ (không giải phương trình).
\loigiai{Theo định lí Viète: $x_1+x_2=6$, $x_1x_2=-16$. Suy ra $x_1^{2}+x_2^{2}=(x_1+x_2)^{2}-2x_1x_2=6^{2}-2\cdot(-16)=68$.}
\end{cauhoi}
\begin{cauhoi}{Bài 26}{9H2-C6-B20-D02-TL-VD-026}
Gọi $x_1$, $x_2$ là hai nghiệm của phương trình $x^{2}+2x=0$. Tính $x_1^{2}+x_2^{2}$ (không giải phương trình).
\loigiai{Theo định lí Viète: $x_1+x_2=-2$, $x_1x_2=0$. Suy ra $x_1^{2}+x_2^{2}=(x_1+x_2)^{2}-2x_1x_2=-2^{2}-2\cdot(0)=4$.}
\end{cauhoi}
\begin{cauhoi}{Bài 27}{9H2-C6-B20-D02-TL-VD-027}
Gọi $x_1$, $x_2$ là hai nghiệm của phương trình $x^{2}-5x-14=0$. Tính $x_1^{2}+x_2^{2}$ (không giải phương trình).
\loigiai{Theo định lí Viète: $x_1+x_2=5$, $x_1x_2=-14$. Suy ra $x_1^{2}+x_2^{2}=(x_1+x_2)^{2}-2x_1x_2=5^{2}-2\cdot(-14)=53$.}
\end{cauhoi}
\begin{cauhoi}{Bài 28}{9H2-C6-B20-D02-TL-VD-028}
Gọi $x_1$, $x_2$ là hai nghiệm của phương trình $x^{2}+2x-8=0$. Tính $x_1^{2}+x_2^{2}$ (không giải phương trình).
\loigiai{Theo định lí Viète: $x_1+x_2=-2$, $x_1x_2=-8$. Suy ra $x_1^{2}+x_2^{2}=(x_1+x_2)^{2}-2x_1x_2=-2^{2}-2\cdot(-8)=20$.}
\end{cauhoi}
\begin{cauhoi}{Bài 29}{9H2-C6-B20-D02-TL-VD-029}
Gọi $x_1$, $x_2$ là hai nghiệm của phương trình $x^{2}+x-12=0$. Tính $x_1^{2}+x_2^{2}$ (không giải phương trình).
\loigiai{Theo định lí Viète: $x_1+x_2=-1$, $x_1x_2=-12$. Suy ra $x_1^{2}+x_2^{2}=(x_1+x_2)^{2}-2x_1x_2=-1^{2}-2\cdot(-12)=25$.}
\end{cauhoi}
\begin{cauhoi}{Bài 30}{9H2-C6-B20-D02-TL-VD-030}
Gọi $x_1$, $x_2$ là hai nghiệm của phương trình $x^{2}+2x-24=0$. Tính $x_1^{2}+x_2^{2}$ (không giải phương trình).
\loigiai{Theo định lí Viète: $x_1+x_2=-2$, $x_1x_2=-24$. Suy ra $x_1^{2}+x_2^{2}=(x_1+x_2)^{2}-2x_1x_2=-2^{2}-2\cdot(-24)=52$.}
\end{cauhoi}
\begin{cauhoi}{Bài 31}{9H2-C6-B20-D02-TL-VD-031}
Gọi $x_1$, $x_2$ là hai nghiệm của phương trình $x^{2}+10x+24=0$. Tính $x_1^{2}+x_2^{2}$ (không giải phương trình).
\loigiai{Theo định lí Viète: $x_1+x_2=-10$, $x_1x_2=24$. Suy ra $x_1^{2}+x_2^{2}=(x_1+x_2)^{2}-2x_1x_2=-10^{2}-2\cdot(24)=52$.}
\end{cauhoi}
\begin{cauhoi}{Bài 32}{9H2-C6-B20-D02-TL-VD-032}
Gọi $x_1$, $x_2$ là hai nghiệm của phương trình $x^{2}-x-6=0$. Tính $x_1^{2}+x_2^{2}$ (không giải phương trình).
\loigiai{Theo định lí Viète: $x_1+x_2=1$, $x_1x_2=-6$. Suy ra $x_1^{2}+x_2^{2}=(x_1+x_2)^{2}-2x_1x_2=1^{2}-2\cdot(-6)=13$.}
\end{cauhoi}
\begin{cauhoi}{Bài 33}{9H2-C6-B20-D02-TL-VD-033}
Gọi $x_1$, $x_2$ là hai nghiệm của phương trình $x^{2}-9x+8=0$. Tính $x_1^{2}+x_2^{2}$ (không giải phương trình).
\loigiai{Theo định lí Viète: $x_1+x_2=9$, $x_1x_2=8$. Suy ra $x_1^{2}+x_2^{2}=(x_1+x_2)^{2}-2x_1x_2=9^{2}-2\cdot(8)=65$.}
\end{cauhoi}
\begin{cauhoi}{Bài 34}{9H2-C6-B20-D02-TL-VD-034}
Gọi $x_1$, $x_2$ là hai nghiệm của phương trình $x^{2}-7x+10=0$. Tính $x_1^{2}+x_2^{2}$ (không giải phương trình).
\loigiai{Theo định lí Viète: $x_1+x_2=7$, $x_1x_2=10$. Suy ra $x_1^{2}+x_2^{2}=(x_1+x_2)^{2}-2x_1x_2=7^{2}-2\cdot(10)=29$.}
\end{cauhoi}
\begin{cauhoi}{Bài 35}{9H2-C6-B20-D02-TL-VD-035}
Gọi $x_1$, $x_2$ là hai nghiệm của phương trình $x^{2}-4x-21=0$. Tính $x_1^{2}+x_2^{2}$ (không giải phương trình).
\loigiai{Theo định lí Viète: $x_1+x_2=4$, $x_1x_2=-21$. Suy ra $x_1^{2}+x_2^{2}=(x_1+x_2)^{2}-2x_1x_2=4^{2}-2\cdot(-21)=58$.}
\end{cauhoi}
\begin{cauhoi}{Bài 36}{9H2-C6-B20-D02-TL-VD-036}
Gọi $x_1$, $x_2$ là hai nghiệm của phương trình $x^{2}-2x-3=0$. Tính $x_1^{2}+x_2^{2}$ (không giải phương trình).
\loigiai{Theo định lí Viète: $x_1+x_2=2$, $x_1x_2=-3$. Suy ra $x_1^{2}+x_2^{2}=(x_1+x_2)^{2}-2x_1x_2=2^{2}-2\cdot(-3)=10$.}
\end{cauhoi}
\begin{cauhoi}{Bài 37}{9H2-C6-B20-D02-TL-VD-037}
Gọi $x_1$, $x_2$ là hai nghiệm của phương trình $x^{2}-10x+16=0$. Tính $x_1^{2}+x_2^{2}$ (không giải phương trình).
\loigiai{Theo định lí Viète: $x_1+x_2=10$, $x_1x_2=16$. Suy ra $x_1^{2}+x_2^{2}=(x_1+x_2)^{2}-2x_1x_2=10^{2}-2\cdot(16)=68$.}
\end{cauhoi}
\begin{cauhoi}{Bài 38}{9H2-C6-B20-D02-TL-VD-038}
Gọi $x_1$, $x_2$ là hai nghiệm của phương trình $x^{2}-x-12=0$. Tính $x_1^{2}+x_2^{2}$ (không giải phương trình).
\loigiai{Theo định lí Viète: $x_1+x_2=1$, $x_1x_2=-12$. Suy ra $x_1^{2}+x_2^{2}=(x_1+x_2)^{2}-2x_1x_2=1^{2}-2\cdot(-12)=25$.}
\end{cauhoi}
\begin{cauhoi}{Bài 39}{9H2-C6-B20-D02-TL-VD-039}
Gọi $x_1$, $x_2$ là hai nghiệm của phương trình $x^{2}+2x-3=0$. Tính $x_1^{2}+x_2^{2}$ (không giải phương trình).
\loigiai{Theo định lí Viète: $x_1+x_2=-2$, $x_1x_2=-3$. Suy ra $x_1^{2}+x_2^{2}=(x_1+x_2)^{2}-2x_1x_2=-2^{2}-2\cdot(-3)=10$.}
\end{cauhoi}
\begin{cauhoi}{Bài 40}{9H2-C6-B20-D02-TL-VD-040}
Gọi $x_1$, $x_2$ là hai nghiệm của phương trình $x^{2}+2x-35=0$. Tính $x_1^{2}+x_2^{2}$ (không giải phương trình).
\loigiai{Theo định lí Viète: $x_1+x_2=-2$, $x_1x_2=-35$. Suy ra $x_1^{2}+x_2^{2}=(x_1+x_2)^{2}-2x_1x_2=-2^{2}-2\cdot(-35)=74$.}
\end{cauhoi}
\begin{cauhoi}{Bài 41}{9H2-C6-B20-D02-TL-VD-041}
Gọi $x_1$, $x_2$ là hai nghiệm của phương trình $x^{2}+6x-16=0$. Tính $x_1^{2}+x_2^{2}$ (không giải phương trình).
\loigiai{Theo định lí Viète: $x_1+x_2=-6$, $x_1x_2=-16$. Suy ra $x_1^{2}+x_2^{2}=(x_1+x_2)^{2}-2x_1x_2=-6^{2}-2\cdot(-16)=68$.}
\end{cauhoi}
\begin{cauhoi}{Bài 42}{9H2-C6-B20-D02-TL-VD-042}
Gọi $x_1$, $x_2$ là hai nghiệm của phương trình $x^{2}+12x+35=0$. Tính $x_1^{2}+x_2^{2}$ (không giải phương trình).
\loigiai{Theo định lí Viète: $x_1+x_2=-12$, $x_1x_2=35$. Suy ra $x_1^{2}+x_2^{2}=(x_1+x_2)^{2}-2x_1x_2=-12^{2}-2\cdot(35)=74$.}
\end{cauhoi}
\begin{cauhoi}{Bài 43}{9H2-C6-B20-D02-TL-VD-043}
Gọi $x_1$, $x_2$ là hai nghiệm của phương trình $x^{2}-6x=0$. Tính $x_1^{2}+x_2^{2}$ (không giải phương trình).
\loigiai{Theo định lí Viète: $x_1+x_2=6$, $x_1x_2=0$. Suy ra $x_1^{2}+x_2^{2}=(x_1+x_2)^{2}-2x_1x_2=6^{2}-2\cdot(0)=36$.}
\end{cauhoi}
\begin{cauhoi}{Bài 44}{9H2-C6-B20-D02-TL-VD-044}
Gọi $x_1$, $x_2$ là hai nghiệm của phương trình $x^{2}-7x+12=0$. Tính $x_1^{2}+x_2^{2}$ (không giải phương trình).
\loigiai{Theo định lí Viète: $x_1+x_2=7$, $x_1x_2=12$. Suy ra $x_1^{2}+x_2^{2}=(x_1+x_2)^{2}-2x_1x_2=7^{2}-2\cdot(12)=25$.}
\end{cauhoi}
\begin{cauhoi}{Bài 45}{9H2-C6-B20-D02-TL-VD-045}
Gọi $x_1$, $x_2$ là hai nghiệm của phương trình $x^{2}-2x=0$. Tính $x_1^{2}+x_2^{2}$ (không giải phương trình).
\loigiai{Theo định lí Viète: $x_1+x_2=2$, $x_1x_2=0$. Suy ra $x_1^{2}+x_2^{2}=(x_1+x_2)^{2}-2x_1x_2=2^{2}-2\cdot(0)=4$.}
\end{cauhoi}
\begin{cauhoi}{Bài 46}{9H2-C6-B20-D02-TL-VD-046}
Gọi $x_1$, $x_2$ là hai nghiệm của phương trình $x^{2}-3x-40=0$. Tính $x_1^{2}+x_2^{2}$ (không giải phương trình).
\loigiai{Theo định lí Viète: $x_1+x_2=3$, $x_1x_2=-40$. Suy ra $x_1^{2}+x_2^{2}=(x_1+x_2)^{2}-2x_1x_2=3^{2}-2\cdot(-40)=89$.}
\end{cauhoi}
\begin{cauhoi}{Bài 47}{9H2-C6-B20-D02-TL-VD-047}
Gọi $x_1$, $x_2$ là hai nghiệm của phương trình $x^{2}+14x+48=0$. Tính $x_1^{2}+x_2^{2}$ (không giải phương trình).
\loigiai{Theo định lí Viète: $x_1+x_2=-14$, $x_1x_2=48$. Suy ra $x_1^{2}+x_2^{2}=(x_1+x_2)^{2}-2x_1x_2=-14^{2}-2\cdot(48)=100$.}
\end{cauhoi}
\begin{cauhoi}{Bài 48}{9H2-C6-B20-D02-TL-VD-048}
Gọi $x_1$, $x_2$ là hai nghiệm của phương trình $x^{2}+5x=0$. Tính $x_1^{2}+x_2^{2}$ (không giải phương trình).
\loigiai{Theo định lí Viète: $x_1+x_2=-5$, $x_1x_2=0$. Suy ra $x_1^{2}+x_2^{2}=(x_1+x_2)^{2}-2x_1x_2=-5^{2}-2\cdot(0)=25$.}
\end{cauhoi}
\begin{cauhoi}{Bài 49}{9H2-C6-B20-D02-TL-VD-049}
Gọi $x_1$, $x_2$ là hai nghiệm của phương trình $x^{2}-x=0$. Tính $x_1^{2}+x_2^{2}$ (không giải phương trình).
\loigiai{Theo định lí Viète: $x_1+x_2=1$, $x_1x_2=0$. Suy ra $x_1^{2}+x_2^{2}=(x_1+x_2)^{2}-2x_1x_2=1^{2}-2\cdot(0)=1$.}
\end{cauhoi}
\begin{cauhoi}{Bài 50}{9H2-C6-B20-D02-TL-VD-050}
Gọi $x_1$, $x_2$ là hai nghiệm của phương trình $x^{2}-4x-32=0$. Tính $x_1^{2}+x_2^{2}$ (không giải phương trình).
\loigiai{Theo định lí Viète: $x_1+x_2=4$, $x_1x_2=-32$. Suy ra $x_1^{2}+x_2^{2}=(x_1+x_2)^{2}-2x_1x_2=4^{2}-2\cdot(-32)=80$.}
\end{cauhoi}
\clearpage\section*{Phần II. Trắc nghiệm}
\begin{cauhoi}{Câu 1}{9H2-C6-B20-D02-TN-VD-001}
Gọi $x_1$, $x_2$ là hai nghiệm của phương trình $x^{2}+6x+8=0$. Tính $x_1^{2}+x_2^{2}$ (không giải phương trình).
\begin{tasks}(2)
\task $20$
\task $36$
\task $52$
\task $21$
\end{tasks}
\dapan{A}
\end{cauhoi}
\begin{cauhoi}{Câu 2}{9H2-C6-B20-D02-TN-VD-002}
Gọi $x_1$, $x_2$ là hai nghiệm của phương trình $x^{2}-4x-5=0$. Tính $x_1^{2}+x_2^{2}$ (không giải phương trình).
\begin{tasks}(2)
\task $26$
\task $6$
\task $16$
\task $27$
\end{tasks}
\dapan{A}
\end{cauhoi}
\begin{cauhoi}{Câu 3}{9H2-C6-B20-D02-TN-VD-003}
Gọi $x_1$, $x_2$ là hai nghiệm của phương trình $x^{2}-x-56=0$. Tính $x_1^{2}+x_2^{2}$ (không giải phương trình).
\begin{tasks}(2)
\task $113$
\task $114$
\task $-111$
\task $1$
\end{tasks}
\dapan{A}
\end{cauhoi}
\begin{cauhoi}{Câu 4}{9H2-C6-B20-D02-TN-VD-004}
Gọi $x_1$, $x_2$ là hai nghiệm của phương trình $x^{2}+9x+8=0$. Tính $x_1^{2}+x_2^{2}$ (không giải phương trình).
\begin{tasks}(2)
\task $66$
\task $65$
\task $97$
\task $81$
\end{tasks}
\dapan{B}
\end{cauhoi}
\begin{cauhoi}{Câu 5}{9H2-C6-B20-D02-TN-VD-005}
Gọi $x_1$, $x_2$ là hai nghiệm của phương trình $x^{2}+7x=0$. Tính $x_1^{2}+x_2^{2}$ (không giải phương trình).
\begin{tasks}(2)
\task $49$
\task $48$
\task $51$
\task $50$
\end{tasks}
\dapan{A}
\end{cauhoi}
\begin{cauhoi}{Câu 6}{9H2-C6-B20-D02-TN-VD-006}
Gọi $x_1$, $x_2$ là hai nghiệm của phương trình $x^{2}-4=0$. Tính $x_1^{2}+x_2^{2}$ (không giải phương trình).
\begin{tasks}(2)
\task $9$
\task $-8$
\task $8$
\task $0$
\end{tasks}
\dapan{C}
\end{cauhoi}
\begin{cauhoi}{Câu 7}{9H2-C6-B20-D02-TN-VD-007}
Gọi $x_1$, $x_2$ là hai nghiệm của phương trình $x^{2}-11x+24=0$. Tính $x_1^{2}+x_2^{2}$ (không giải phương trình).
\begin{tasks}(2)
\task $169$
\task $121$
\task $74$
\task $73$
\end{tasks}
\dapan{D}
\end{cauhoi}
\begin{cauhoi}{Câu 8}{9H2-C6-B20-D02-TN-VD-008}
Gọi $x_1$, $x_2$ là hai nghiệm của phương trình $x^{2}-4x+3=0$. Tính $x_1^{2}+x_2^{2}$ (không giải phương trình).
\begin{tasks}(2)
\task $10$
\task $11$
\task $16$
\task $22$
\end{tasks}
\dapan{A}
\end{cauhoi}
\begin{cauhoi}{Câu 9}{9H2-C6-B20-D02-TN-VD-009}
Gọi $x_1$, $x_2$ là hai nghiệm của phương trình $x^{2}-25=0$. Tính $x_1^{2}+x_2^{2}$ (không giải phương trình).
\begin{tasks}(2)
\task $-50$
\task $0$
\task $51$
\task $50$
\end{tasks}
\dapan{D}
\end{cauhoi}
\begin{cauhoi}{Câu 10}{9H2-C6-B20-D02-TN-VD-010}
Gọi $x_1$, $x_2$ là hai nghiệm của phương trình $x^{2}-8x=0$. Tính $x_1^{2}+x_2^{2}$ (không giải phương trình).
\begin{tasks}(2)
\task $65$
\task $63$
\task $66$
\task $64$
\end{tasks}
\dapan{D}
\end{cauhoi}
\begin{cauhoi}{Câu 11}{9H2-C6-B20-D02-TN-VD-011}
Gọi $x_1$, $x_2$ là hai nghiệm của phương trình $x^{2}+4x-21=0$. Tính $x_1^{2}+x_2^{2}$ (không giải phương trình).
\begin{tasks}(2)
\task $58$
\task $59$
\task $-26$
\task $16$
\end{tasks}
\dapan{A}
\end{cauhoi}
\begin{cauhoi}{Câu 12}{9H2-C6-B20-D02-TN-VD-012}
Gọi $x_1$, $x_2$ là hai nghiệm của phương trình $x^{2}-3x-4=0$. Tính $x_1^{2}+x_2^{2}$ (không giải phương trình).
\begin{tasks}(2)
\task $17$
\task $1$
\task $18$
\task $9$
\end{tasks}
\dapan{A}
\end{cauhoi}
\begin{cauhoi}{Câu 13}{9H2-C6-B20-D02-TN-VD-013}
Gọi $x_1$, $x_2$ là hai nghiệm của phương trình $x^{2}-12x+35=0$. Tính $x_1^{2}+x_2^{2}$ (không giải phương trình).
\begin{tasks}(2)
\task $214$
\task $75$
\task $74$
\task $144$
\end{tasks}
\dapan{C}
\end{cauhoi}
\begin{cauhoi}{Câu 14}{9H2-C6-B20-D02-TN-VD-014}
Gọi $x_1$, $x_2$ là hai nghiệm của phương trình $x^{2}+8x=0$. Tính $x_1^{2}+x_2^{2}$ (không giải phương trình).
\begin{tasks}(2)
\task $66$
\task $64$
\task $65$
\task $63$
\end{tasks}
\dapan{B}
\end{cauhoi}
\begin{cauhoi}{Câu 15}{9H2-C6-B20-D02-TN-VD-015}
Gọi $x_1$, $x_2$ là hai nghiệm của phương trình $x^{2}+x=0$. Tính $x_1^{2}+x_2^{2}$ (không giải phương trình).
\begin{tasks}(2)
\task $3$
\task $2$
\task $0$
\task $1$
\end{tasks}
\dapan{D}
\end{cauhoi}
\begin{cauhoi}{Câu 16}{9H2-C6-B20-D02-TN-VD-016}
Gọi $x_1$, $x_2$ là hai nghiệm của phương trình $x^{2}+x-42=0$. Tính $x_1^{2}+x_2^{2}$ (không giải phương trình).
\begin{tasks}(2)
\task $-83$
\task $85$
\task $1$
\task $86$
\end{tasks}
\dapan{B}
\end{cauhoi}
\begin{cauhoi}{Câu 17}{9H2-C6-B20-D02-TN-VD-017}
Gọi $x_1$, $x_2$ là hai nghiệm của phương trình $x^{2}+2x-48=0$. Tính $x_1^{2}+x_2^{2}$ (không giải phương trình).
\begin{tasks}(2)
\task $-92$
\task $101$
\task $4$
\task $100$
\end{tasks}
\dapan{D}
\end{cauhoi}
\begin{cauhoi}{Câu 18}{9H2-C6-B20-D02-TN-VD-018}
Gọi $x_1$, $x_2$ là hai nghiệm của phương trình $x^{2}-x-30=0$. Tính $x_1^{2}+x_2^{2}$ (không giải phương trình).
\begin{tasks}(2)
\task $61$
\task $-59$
\task $1$
\task $62$
\end{tasks}
\dapan{A}
\end{cauhoi}
\begin{cauhoi}{Câu 19}{9H2-C6-B20-D02-TN-VD-019}
Gọi $x_1$, $x_2$ là hai nghiệm của phương trình $x^{2}+x-20=0$. Tính $x_1^{2}+x_2^{2}$ (không giải phương trình).
\begin{tasks}(2)
\task $1$
\task $41$
\task $-39$
\task $42$
\end{tasks}
\dapan{B}
\end{cauhoi}
\begin{cauhoi}{Câu 20}{9H2-C6-B20-D02-TN-VD-020}
Gọi $x_1$, $x_2$ là hai nghiệm của phương trình $x^{2}-2x-48=0$. Tính $x_1^{2}+x_2^{2}$ (không giải phương trình).
\begin{tasks}(2)
\task $4$
\task $-92$
\task $100$
\task $101$
\end{tasks}
\dapan{C}
\end{cauhoi}
\begin{cauhoi}{Câu 21}{9H2-C6-B20-D02-TN-VD-021}
Gọi $x_1$, $x_2$ là hai nghiệm của phương trình $x^{2}+13x+42=0$. Tính $x_1^{2}+x_2^{2}$ (không giải phương trình).
\begin{tasks}(2)
\task $169$
\task $85$
\task $253$
\task $86$
\end{tasks}
\dapan{B}
\end{cauhoi}
\begin{cauhoi}{Câu 22}{9H2-C6-B20-D02-TN-VD-022}
Gọi $x_1$, $x_2$ là hai nghiệm của phương trình $x^{2}+4x-32=0$. Tính $x_1^{2}+x_2^{2}$ (không giải phương trình).
\begin{tasks}(2)
\task $16$
\task $-48$
\task $81$
\task $80$
\end{tasks}
\dapan{D}
\end{cauhoi}
\begin{cauhoi}{Câu 23}{9H2-C6-B20-D02-TN-VD-023}
Gọi $x_1$, $x_2$ là hai nghiệm của phương trình $x^{2}+4x-5=0$. Tính $x_1^{2}+x_2^{2}$ (không giải phương trình).
\begin{tasks}(2)
\task $27$
\task $26$
\task $16$
\task $6$
\end{tasks}
\dapan{B}
\end{cauhoi}
\begin{cauhoi}{Câu 24}{9H2-C6-B20-D02-TN-VD-024}
Gọi $x_1$, $x_2$ là hai nghiệm của phương trình $x^{2}-3x-18=0$. Tính $x_1^{2}+x_2^{2}$ (không giải phương trình).
\begin{tasks}(2)
\task $9$
\task $45$
\task $46$
\task $-27$
\end{tasks}
\dapan{B}
\end{cauhoi}
\begin{cauhoi}{Câu 25}{9H2-C6-B20-D02-TN-VD-025}
Gọi $x_1$, $x_2$ là hai nghiệm của phương trình $x^{2}-6x-16=0$. Tính $x_1^{2}+x_2^{2}$ (không giải phương trình).
\begin{tasks}(2)
\task $4$
\task $36$
\task $68$
\task $69$
\end{tasks}
\dapan{C}
\end{cauhoi}
\begin{cauhoi}{Câu 26}{9H2-C6-B20-D02-TN-VD-026}
Gọi $x_1$, $x_2$ là hai nghiệm của phương trình $x^{2}+2x=0$. Tính $x_1^{2}+x_2^{2}$ (không giải phương trình).
\begin{tasks}(2)
\task $5$
\task $6$
\task $4$
\task $3$
\end{tasks}
\dapan{C}
\end{cauhoi}
\begin{cauhoi}{Câu 27}{9H2-C6-B20-D02-TN-VD-027}
Gọi $x_1$, $x_2$ là hai nghiệm của phương trình $x^{2}-5x-14=0$. Tính $x_1^{2}+x_2^{2}$ (không giải phương trình).
\begin{tasks}(2)
\task $53$
\task $54$
\task $-3$
\task $25$
\end{tasks}
\dapan{A}
\end{cauhoi}
\begin{cauhoi}{Câu 28}{9H2-C6-B20-D02-TN-VD-028}
Gọi $x_1$, $x_2$ là hai nghiệm của phương trình $x^{2}+2x-8=0$. Tính $x_1^{2}+x_2^{2}$ (không giải phương trình).
\begin{tasks}(2)
\task $-12$
\task $4$
\task $20$
\task $21$
\end{tasks}
\dapan{C}
\end{cauhoi}
\begin{cauhoi}{Câu 29}{9H2-C6-B20-D02-TN-VD-029}
Gọi $x_1$, $x_2$ là hai nghiệm của phương trình $x^{2}+x-12=0$. Tính $x_1^{2}+x_2^{2}$ (không giải phương trình).
\begin{tasks}(2)
\task $-23$
\task $1$
\task $26$
\task $25$
\end{tasks}
\dapan{D}
\end{cauhoi}
\begin{cauhoi}{Câu 30}{9H2-C6-B20-D02-TN-VD-030}
Gọi $x_1$, $x_2$ là hai nghiệm của phương trình $x^{2}+2x-24=0$. Tính $x_1^{2}+x_2^{2}$ (không giải phương trình).
\begin{tasks}(2)
\task $4$
\task $53$
\task $52$
\task $-44$
\end{tasks}
\dapan{C}
\end{cauhoi}
\begin{cauhoi}{Câu 31}{9H2-C6-B20-D02-TN-VD-031}
Gọi $x_1$, $x_2$ là hai nghiệm của phương trình $x^{2}+10x+24=0$. Tính $x_1^{2}+x_2^{2}$ (không giải phương trình).
\begin{tasks}(2)
\task $52$
\task $148$
\task $53$
\task $100$
\end{tasks}
\dapan{A}
\end{cauhoi}
\begin{cauhoi}{Câu 32}{9H2-C6-B20-D02-TN-VD-032}
Gọi $x_1$, $x_2$ là hai nghiệm của phương trình $x^{2}-x-6=0$. Tính $x_1^{2}+x_2^{2}$ (không giải phương trình).
\begin{tasks}(2)
\task $1$
\task $13$
\task $-11$
\task $14$
\end{tasks}
\dapan{B}
\end{cauhoi}
\begin{cauhoi}{Câu 33}{9H2-C6-B20-D02-TN-VD-033}
Gọi $x_1$, $x_2$ là hai nghiệm của phương trình $x^{2}-9x+8=0$. Tính $x_1^{2}+x_2^{2}$ (không giải phương trình).
\begin{tasks}(2)
\task $66$
\task $81$
\task $97$
\task $65$
\end{tasks}
\dapan{D}
\end{cauhoi}
\begin{cauhoi}{Câu 34}{9H2-C6-B20-D02-TN-VD-034}
Gọi $x_1$, $x_2$ là hai nghiệm của phương trình $x^{2}-7x+10=0$. Tính $x_1^{2}+x_2^{2}$ (không giải phương trình).
\begin{tasks}(2)
\task $29$
\task $49$
\task $30$
\task $69$
\end{tasks}
\dapan{A}
\end{cauhoi}
\begin{cauhoi}{Câu 35}{9H2-C6-B20-D02-TN-VD-035}
Gọi $x_1$, $x_2$ là hai nghiệm của phương trình $x^{2}-4x-21=0$. Tính $x_1^{2}+x_2^{2}$ (không giải phương trình).
\begin{tasks}(2)
\task $58$
\task $-26$
\task $59$
\task $16$
\end{tasks}
\dapan{A}
\end{cauhoi}
\begin{cauhoi}{Câu 36}{9H2-C6-B20-D02-TN-VD-036}
Gọi $x_1$, $x_2$ là hai nghiệm của phương trình $x^{2}-2x-3=0$. Tính $x_1^{2}+x_2^{2}$ (không giải phương trình).
\begin{tasks}(2)
\task $11$
\task $4$
\task $-2$
\task $10$
\end{tasks}
\dapan{D}
\end{cauhoi}
\begin{cauhoi}{Câu 37}{9H2-C6-B20-D02-TN-VD-037}
Gọi $x_1$, $x_2$ là hai nghiệm của phương trình $x^{2}-10x+16=0$. Tính $x_1^{2}+x_2^{2}$ (không giải phương trình).
\begin{tasks}(2)
\task $100$
\task $68$
\task $69$
\task $132$
\end{tasks}
\dapan{B}
\end{cauhoi}
\begin{cauhoi}{Câu 38}{9H2-C6-B20-D02-TN-VD-038}
Gọi $x_1$, $x_2$ là hai nghiệm của phương trình $x^{2}-x-12=0$. Tính $x_1^{2}+x_2^{2}$ (không giải phương trình).
\begin{tasks}(2)
\task $26$
\task $-23$
\task $1$
\task $25$
\end{tasks}
\dapan{D}
\end{cauhoi}
\begin{cauhoi}{Câu 39}{9H2-C6-B20-D02-TN-VD-039}
Gọi $x_1$, $x_2$ là hai nghiệm của phương trình $x^{2}+2x-3=0$. Tính $x_1^{2}+x_2^{2}$ (không giải phương trình).
\begin{tasks}(2)
\task $4$
\task $11$
\task $10$
\task $-2$
\end{tasks}
\dapan{C}
\end{cauhoi}
\begin{cauhoi}{Câu 40}{9H2-C6-B20-D02-TN-VD-040}
Gọi $x_1$, $x_2$ là hai nghiệm của phương trình $x^{2}+2x-35=0$. Tính $x_1^{2}+x_2^{2}$ (không giải phương trình).
\begin{tasks}(2)
\task $4$
\task $75$
\task $-66$
\task $74$
\end{tasks}
\dapan{D}
\end{cauhoi}
\begin{cauhoi}{Câu 41}{9H2-C6-B20-D02-TN-VD-041}
Gọi $x_1$, $x_2$ là hai nghiệm của phương trình $x^{2}+6x-16=0$. Tính $x_1^{2}+x_2^{2}$ (không giải phương trình).
\begin{tasks}(2)
\task $68$
\task $69$
\task $4$
\task $36$
\end{tasks}
\dapan{A}
\end{cauhoi}
\begin{cauhoi}{Câu 42}{9H2-C6-B20-D02-TN-VD-042}
Gọi $x_1$, $x_2$ là hai nghiệm của phương trình $x^{2}+12x+35=0$. Tính $x_1^{2}+x_2^{2}$ (không giải phương trình).
\begin{tasks}(2)
\task $75$
\task $144$
\task $74$
\task $214$
\end{tasks}
\dapan{C}
\end{cauhoi}
\begin{cauhoi}{Câu 43}{9H2-C6-B20-D02-TN-VD-043}
Gọi $x_1$, $x_2$ là hai nghiệm của phương trình $x^{2}-6x=0$. Tính $x_1^{2}+x_2^{2}$ (không giải phương trình).
\begin{tasks}(2)
\task $37$
\task $35$
\task $38$
\task $36$
\end{tasks}
\dapan{D}
\end{cauhoi}
\begin{cauhoi}{Câu 44}{9H2-C6-B20-D02-TN-VD-044}
Gọi $x_1$, $x_2$ là hai nghiệm của phương trình $x^{2}-7x+12=0$. Tính $x_1^{2}+x_2^{2}$ (không giải phương trình).
\begin{tasks}(2)
\task $49$
\task $25$
\task $26$
\task $73$
\end{tasks}
\dapan{B}
\end{cauhoi}
\begin{cauhoi}{Câu 45}{9H2-C6-B20-D02-TN-VD-045}
Gọi $x_1$, $x_2$ là hai nghiệm của phương trình $x^{2}-2x=0$. Tính $x_1^{2}+x_2^{2}$ (không giải phương trình).
\begin{tasks}(2)
\task $4$
\task $6$
\task $3$
\task $5$
\end{tasks}
\dapan{A}
\end{cauhoi}
\begin{cauhoi}{Câu 46}{9H2-C6-B20-D02-TN-VD-046}
Gọi $x_1$, $x_2$ là hai nghiệm của phương trình $x^{2}-3x-40=0$. Tính $x_1^{2}+x_2^{2}$ (không giải phương trình).
\begin{tasks}(2)
\task $9$
\task $89$
\task $-71$
\task $90$
\end{tasks}
\dapan{B}
\end{cauhoi}
\begin{cauhoi}{Câu 47}{9H2-C6-B20-D02-TN-VD-047}
Gọi $x_1$, $x_2$ là hai nghiệm của phương trình $x^{2}+14x+48=0$. Tính $x_1^{2}+x_2^{2}$ (không giải phương trình).
\begin{tasks}(2)
\task $101$
\task $100$
\task $196$
\task $292$
\end{tasks}
\dapan{B}
\end{cauhoi}
\begin{cauhoi}{Câu 48}{9H2-C6-B20-D02-TN-VD-048}
Gọi $x_1$, $x_2$ là hai nghiệm của phương trình $x^{2}+5x=0$. Tính $x_1^{2}+x_2^{2}$ (không giải phương trình).
\begin{tasks}(2)
\task $24$
\task $27$
\task $25$
\task $26$
\end{tasks}
\dapan{C}
\end{cauhoi}
\begin{cauhoi}{Câu 49}{9H2-C6-B20-D02-TN-VD-049}
Gọi $x_1$, $x_2$ là hai nghiệm của phương trình $x^{2}-x=0$. Tính $x_1^{2}+x_2^{2}$ (không giải phương trình).
\begin{tasks}(2)
\task $0$
\task $2$
\task $3$
\task $1$
\end{tasks}
\dapan{D}
\end{cauhoi}
\begin{cauhoi}{Câu 50}{9H2-C6-B20-D02-TN-VD-050}
Gọi $x_1$, $x_2$ là hai nghiệm của phương trình $x^{2}-4x-32=0$. Tính $x_1^{2}+x_2^{2}$ (không giải phương trình).
\begin{tasks}(2)
\task $80$
\task $-48$
\task $81$
\task $16$
\end{tasks}
\dapan{A}
\end{cauhoi}
\clearpage\section*{Bảng đáp án trắc nghiệm}
\begin{center}\begin{tabular}{|c|c|c|c|c|c|c|c|c|c|}
\hline
1. A & 2. A & 3. A & 4. B & 5. A & 6. C & 7. D & 8. A & 9. D & 10. D\\\hline
11. A & 12. A & 13. C & 14. B & 15. D & 16. B & 17. D & 18. A & 19. B & 20. C\\\hline
21. B & 22. D & 23. B & 24. B & 25. C & 26. C & 27. A & 28. C & 29. D & 30. C\\\hline
31. A & 32. B & 33. D & 34. A & 35. A & 36. D & 37. B & 38. D & 39. C & 40. D\\\hline
41. A & 42. C & 43. D & 44. B & 45. A & 46. B & 47. B & 48. C & 49. D & 50. A\\\hline
\end{tabular}\end{center}
\end{document}
